\chapter{Semantic Caching \& Discretization}
\label{ch:semantic_caching}

Even with the implementation of Singleflight request coalescing (Chapter 8), executing full LLM inference for 400 agents at 1,000 Ticks Per Second remains physically impossible on consumer hardware. To achieve true high-frequency scale, the simulator must learn to bypass the GPU entirely for the vast majority of agent decisions. This is accomplished through advanced Structured Intermediate Representation Caching (SIRC).

\section{The Concept of Semantic Caching}

In traditional software, a cache stores the exact output of a deterministic function given a specific input string. If the exact input is seen again, the cache returns the output in $O(1)$ time. 

However, LLMs deal with unstructured natural language and continuous, high-dimensional state vectors (such as the LOB state). Two market states are never byte-for-byte identical, but they may be \textit{semantically identical} in their financial implications. 

Semantic caching intercepts the agent's input, converts it into a mathematical vector (an embedding), and searches a historical database for vectors that are "close enough." If a highly similar historical state is found, the agent bypasses the GPU and instantly executes the cached financial decision.

\section{Matryoshka Representation Learning (MRL)}

Converting a complex LOB state and a narrative shock into a vector typically requires passing it through a dense embedding model (e.g., OpenAI's `text-embedding-3`). These models output massive floating-point vectors, often 1024 or 1536 dimensions wide. 

Storing millions of 1024-dimensional, 32-bit floating-point (FP32) vectors in memory quickly exhausts RAM. Furthermore, computing the Cosine Similarity between thousands of these vectors takes far too long for a microsecond-latency event loop.

To compress this data, AMPS utilizes \textbf{Matryoshka Representation Learning (MRL)} \cite{kusupati2022matryoshka}. Traditional embedding models distribute semantic meaning evenly across all 1024 dimensions. MRL, however, forces the neural network to encode the most critical, broad semantic meaning into the very first few dimensions, with subsequent dimensions adding only fine-grained nuances (like a Russian Matryoshka nesting doll).

This allows the AMPS engine to safely truncate the 1024D vector down to its first 128 dimensions. This aggressive truncation preserves up to 85\% of the topological semantics while instantly reducing the memory and computational footprint by an order of magnitude.

\section{1-Bit Binary Quantization (BQ)}

While 128D FP32 vectors are smaller, calculating floating-point mathematics is still fundamentally slower than integer or bitwise arithmetic. To achieve absolute maximum search throughput, AMPS applies \textbf{1-Bit Binary Quantization (BQ)} to the truncated MRL embeddings.

Binary quantization converts the floating-point vector into a string of 1s and 0s. Mathematically, if a dimension's value is greater than zero, it becomes a 1; if less than or equal to zero, it becomes a 0.

\begin{equation}
    b_i = \begin{cases} 
      1 & x_i > 0 \\
      0 & x_i \leq 0 
   \end{cases}
\end{equation}

By applying this thresholding, an entire agent's cognitive and financial state is compressed from a massive matrix down to a single 128-bit array (exactly 16 bytes of memory).

However, in a 128-bit Hamming space, this extreme quantization noise introduces a $\sim 10$-$15\%$ angular estimation margin of error, creating a high risk of false-positive cache hits where subtle but critical market shifts are incorrectly matched. To eliminate these false positives, AMPS employs a \textbf{Secondary FP32 Verification Step}. The 1-bit Hamming search is used strictly as an ultra-fast first-pass filter to retrieve the top-k candidates, which are then re-ranked using their original dense FP32 representations before executing the cache hit.

\begin{figure}[htbp]
\centering
\begin{tikzpicture}[
    scale=0.9, transform shape,
    box/.style={draw, rectangle, minimum width=3cm, align=center, minimum height=1cm},
    arr/.style={->, thick, >=stealth}
]
    \node[box, fill=blue!15] (dense) {Dense FP32 Embed\\(1024 Dimensions\\$\approx 4$ Kilobytes)};
    \node[box, fill=green!15] (trunc) [right=1cm of dense] {MRL Truncation\\(128 Dimensions\\$\approx 512$ Bytes)};
    \node[box, fill=red!15] (quant) [right=1cm of trunc] {1-Bit Quantization\\(128 bits = 16 Bytes)};
    
    \draw[arr] (dense) -- (trunc);
    \draw[arr] (trunc) -- (quant);
\end{tikzpicture}
\caption{The semantic compression pipeline. Massive text and financial arrays are funneled down into a hyper-dense 16-byte binary signature.}
\label{fig:mrl_quant_detail}
\end{figure}

\section{Ternary Weight Quantization (1.58-bit)}

Beyond caching, the inference engine itself must be heavily compressed to fit within the memory constraints of a single workstation. Standard LLMs utilize 16-bit floating-point (FP16) parameters. 

AMPS employs ultra-low-bitwidth ternary quantization \cite{ma2024bitnet}. By forcing the massive weight matrices of the neural network to assume only three discrete values ($-1, 0, 1$), the model consumes approximately 1.58 bits per parameter. This aggressive compression drastically reduces VRAM requirements for model weights. 

Maintaining individual Key-Value (KV) caches for 400 agent context windows requires approximately 40 GB of total memory. By utilizing an NVIDIA RTX 5070 Ti equipped with 16GB of VRAM, AMPS cannot pin all agent KV-caches directly in hardware memory simultaneously. 

To resolve PCIe bus contention while respecting this ceiling, AMPS implements a Tiered KV-Cache Pinning Policy (detailed in Appendix \ref{ch:hardware_optimization_blueprint}). The top 50-80 high-frequency liquidity providers and institutional whales are statically pinned directly in dedicated GPU VRAM ($\sim$1.8GB to 2.4GB total allocation), achieving sub-500 nanosecond memory access. The remaining 320-350 low-frequency retail agents have their dormant states paged into pinned host DDR5 RAM (\texttt{torch.cuda.HostRegister}) and streamed asynchronously over PCIe Gen 5 only when active trading is triggered. This tiered hierarchy eliminates over 80\% of bus traffic during standard simulation ticks, ensuring that the Gracemont E-core matching engine maintains isolated, uncontested access to system memory bandwidth.
